\section{Semigroups}\label{sec:semigroups}

We list here several basic algebraic structures that we will use mostly as building blocks for more complicated structures.

\paragraph{Operations}

The \hyperref[def:logical_theory]{logical theories} presented here are \hyperref[def:fol_equational_theory]{equational}. Their signatures have only \hyperref[def:fol_signature/function]{function symbols}, so their \hyperref[def:fol_semantics/model]{models} are determined by a set with \hyperref[def:operation_on_set]{operations}. As shown in \cref{thm:fol_equational_theory_models}, these structures are particularly well-behaved and even have canonical \hyperref[def:fol_free_term_model]{term models}.

As discussed in \cref{rem:mathematical_logic_conventions/structure_pairs}, following popular convention, we will identify these models with their universes, with the operations associated implicitly.

\begin{figure}[!ht]
  \begin{center}
    \begin{forest}
      for tree={grow=north}
      [
        {\hyperref[def:group]{Group}}
          [
            {\hyperref[def:set_with_involution]{Set with involution}}, edge label={node[midway, right=1em]{\( x^{-1} \) (unary)}}
          ]
          [
            {\hyperref[def:monoid]{Monoid}},
              [{\hyperref[def:semigroup]{Semigroup}}, edge label={node[midway, right=1em]{\( x \cdot y \) (binary)}}]
              [{\hyperref[def:pointed_set]{Pointed set}}, edge label={node[midway, left=1em]{\( e \) (nullary)}}]
          ]
      ]
    \end{forest}
  \end{center}
  \caption{Generalizations of groups}\label{fig:monoid_hierarchy}
\end{figure}

\begin{definition}\label{def:operation_on_set}\mcite[132]{Birkhoff1967LatticeTheory}
  An \term[bg=операция (\cite[def. IV.1]{ГеновМиховскиМоллов1991Алгебра}), ru=опярация (\cite[def. 1.4]{Эдельман1975Логика})]{operation} of arity \( n \) on a \hyperref[def:set]{plain set} \( A \) is a \hyperref[def:function]{function} from \( A^n \) to \( A \).
\end{definition}
\begin{comments}
  \item The arity of an operation is not to be confused with the corresponding relation arity as defined in \cref{def:relation/arity} --- functions are always binary relations.
\end{comments}

\begin{definition}\label{def:operation_arity_terminology}\mcite[132]{Birkhoff1967LatticeTheory}
  We introduce the following terminology based on the arity \( n \) of an \hyperref[def:operation_on_set]{operation}:
  \begin{itemize}
    \item \term[ru=нульарная (операция) (\cite[108]{Курош1973ОбщаяАлгебра})]{Nullary} for \( n = 0 \).
    \item \term[bg=унарна (операция) (\cite[75]{ГеновМиховскиМоллов1991Алгебра}), ru=унарная (операция) (\cite[def. 1.4]{Эдельман1975Логика})]{Unary} for \( n = 1 \).
    \item \term[bg=бинарна (операция) (\cite[def. IV.1]{ГеновМиховскиМоллов1991Алгебра}), ru=бинарная (операция) (\cite[def. 1.4]{Эдельман1975Логика})]{Binary} for \( n = 2 \).
    \item \term[bg=тернарна (операция) (\cite[75]{ГеновМиховскиМоллов1991Алгебра}), ru=тернарная (операция) (\cite[def. 1.4]{Эдельман1975Логика})]{Ternary} for \( n = 3 \).
  \end{itemize}
\end{definition}

\paragraph{Pointed sets}

\begin{definition}\label{def:pointed_set}\mcite[26]{MacLane1998CategoryTheory}
  A \term[ru=множество с отмеченной точкой (\cite[16]{ЦаленкоШульгейфер1974ОсновыТеорииКатегорий})]{pointed set} is a simple algebraic structure --- a nonempty set \( X \) with a distinguished element \( e \), called the \term{base point} of \( X \).

  Pointed sets have the following metamathematical properties:
  \begin{thmenum}
    \thmitem{def:pointed_set/theory}\mimprovised Pointed sets can be described by an empty \hyperref[def:fol_theory]{first-order theory} over a \hyperref[def:fol_signature]{signature} with one nullary \hyperref[def:fol_signature/function]{function symbol}, \( \syn1 \).

    This theory is vacuously \hyperref[def:fol_equational_theory]{equational}. As shown in \cref{thm:fol_equational_theory_models}, models of equational theories are particularly well-behaved.

    \thmitem{def:pointed_set/morphism} Given pointed sets \( (X, e_X) \) and \( (Y, e_Y) \), we say that the function \( \varphi: X \to Y \) is a \term[en=map (between pointed sets) (\cite[26]{MacLane1998CategoryTheory})]{pointed set homomorphism} if
    \begin{equation}\label{eq:def:pointed_set/morphism}
      \varphi(e_X) = e_Y.
    \end{equation}

    These are precisely the \hyperref[def:fol_homomorphism]{first-order homomorphisms}, which makes the concepts from \fullref{sec:first_order_structures} applicable to them.

    \thmitem{def:pointed_set/category} For a fixed \hyperref[def:grothendieck_universe]{Grothendieck universe} \( \mscrU \), the first-order theory of pointed sets naturally gives rise to a \hyperref[def:fol_theory_model_functor/obj]{category of \( \mscrU \)-small models}.

    We denote this category by \( \ucat{Set}_* \).

    \thmitem{def:pointed_set/subobject}\mimprovised A \term{pointed subset} of \( (X, e) \) is a subset containing \( e \).

    It follows from \cref{thm:fol_equational_theory_models/substructure} that the pointed subsets of \( (X, e) \) are precisely its \hyperref[def:fol_substructure/submodel]{first-order submodels}, and from \cref{thm:fol_categorical_subobjects} that they are its \hyperref[def:categorical_subobject]{categorical subobjects}.

    Furthermore, as a consequence of \cref{thm:fol_equational_theory_models/image}, the image of a pointed set homomorphism is a pointed subset of its codomain.

    \thmitem{def:pointed_set/trivial}\mimprovised In accordance with \cref{def:trivial_object}, we will call a one-element pointed set \term{trivial} because, as shown in \cref{thm:equational_theory_zero_morphisms/trivial}, it is a zero object in \( \ucat{Set}_* \).
  \end{thmenum}
\end{definition}

\paragraph{Sets with involutions}

\begin{definition}\label{def:involution}\mcite[36]{Schechter1997AnalysisHandbook}
  We say that an \hyperref[def:function/endofunction]{endofunction} \( f: A \to A \) is an \term{involution} if \( f \) is its own \hyperref[def:function_invertibility/bijective]{inverse}.
\end{definition}
\begin{comments}
  \item Involutions generalize to categories --- see \cref{def:morphism_invertibility/involution} --- which makes the term slightly ambiguous.
\end{comments}

\begin{definition}\label{def:set_with_involution}\mimprovised
  A \term{set with an involution} is, unsurprisingly, a \hyperref[def:set]{set} \( X \) equipped with a unary \hyperref[def:involution]{involution} \( \oline{\anon} \).

  Sets with involutions have the following metamathematical properties:
  \begin{thmenum}
    \thmitem{def:set_with_involution/theory} We can describe involutions using \hyperref[def:first_order_logic]{first-order logic}. Over the \hyperref[def:fol_signature]{signature} with one \hyperref[def:fol_signature/notation]{condensed} unary \hyperref[def:fol_signature/function]{function symbol}, \( \syninv \), an \hyperref[def:fol_equational_theory]{equational theory} of sets with involutions is generated by the following axiom:
    \begin{equation}\label{eq:def:set_with_involution/theory/axiom}
      \qforall \synx (\syninv \syninv \synx \syneq \synx).
    \end{equation}

    There is one importance nuance --- we generally allow empty sets with involution, however we generally disallow empty \hyperref[def:fol_semantics/model]{empty first-order models}. We discuss this issue in \cref{rem:fol_empty_universe/semantics}.

    \thmitem{def:set_with_involution/morphism} Given sets \( X \) and \( Y \) with involutions, we say that the function \( \varphi: X \to Y \) is a \term{homomorphism} if
    \begin{equation}\label{eq:def:set_with_involution/morphism}
      \varphi(\oline x) = \oline{\varphi(x)}.
    \end{equation}

    As long as we allow empty universes for first-order models, these are precisely the \hyperref[def:fol_homomorphism]{first-order homomorphisms}, which makes the concepts from \fullref{sec:first_order_structures} applicable to them.

    \thmitem{def:set_with_involution/category} For a fixed \hyperref[def:grothendieck_universe]{Grothendieck universe} \( \mscrU \), we can consider the family of all \( \mscrU \)-small sets with involution along with their homomorphisms.

    This generalizes the \hyperref[def:fol_theory_model_functor/obj]{category of \( \mscrU \)-small first-order models} by also allowing empty sets.

    \thmitem{def:set_with_involution/subobject} A \term{subset with involution} of \( X \) is any subset of \( X \) endowed with the same operation.

    As long as we allow empty universes for first-order models, \cref{thm:fol_equational_theory_models/substructure} implies that the subsets with involution of \( X \) are precisely its \hyperref[def:fol_substructure/submodel]{first-order submodels}, and from \cref{thm:fol_categorical_subobjects} that they are its \hyperref[def:categorical_subobject]{categorical subobjects}.

    Furthermore, as a consequence of \cref{thm:fol_equational_theory_models/image}, the \hyperref[def:set_valued_map/image]{image} of a homomorphism \( \varphi: X \to Y \) is a subset of \( Y \).
  \end{thmenum}
\end{definition}

\paragraph{Binary operations}

\begin{definition}\label{def:binary_operation}
  \hyperref[def:operation_arity_terminology]{Binary} \hyperref[def:operation_on_set]{operations} are of special importance.

  Let \( {\cdot} \) be a an \hyperref[def:fol_signature/notation]{infix} binary \hyperref[def:fol_signature/function]{function symbol}. The following \hyperref[def:fol_formula]{formulas}, or rather their \hyperref[def:fol_quantifier_closure]{universal closures}, are commonly used as axioms:
  \begin{thmenum}
    \thmitem{def:binary_operation/associative}\mcite[\S 1.18(\( \alpha \))]{Ляпин1960Полугруппы} \term[bg=асоциативен закон (\cite[11]{ГеновМиховскиМоллов1991Алгебра}), ru=ассоциативное действие, en=associativity (\cite[12]{GowersEtAl2008PrincetonCompanion})]{Associativity}:
    \begin{equation}\label{eq:def:binary_operation/associative}
      (\synx \cdot \syny) \cdot \synz \syneq \synx \cdot (\syny \cdot \synz).
    \end{equation}

    \thmitem{def:binary_operation/commutative}\mcite[\S 1.18(\( \beta \))]{Ляпин1960Полугруппы} \term[bg=комутативен (закон) (\cite[11]{ГеновМиховскиМоллов1991Алгебра}), ru=коммутативное (действие), en=commutative (binary operation) (\cite[12]{GowersEtAl2008PrincetonCompanion})]{Commutativity}:
    \begin{equation}\label{eq:def:binary_operation/commutative}
      \synx \cdot \syny \syneq \syny \cdot \synx.
    \end{equation}

    \thmitem{def:binary_operation/idempotent}\mcite[114]{БелоусовТкачёв2004ДискретнаяМатематика} \term[ru=идемпотентная (операция), en=idempotent (operation) (\cite[\S A.2.1]{Mimram2020ProgramEqualsProof})]{Idempotence}:
    \begin{equation}\label{eq:def:binary_operation/idempotent}
      \synx \cdot \synx \syneq \synx.
    \end{equation}

    \thmitem{def:binary_operation/cancellative}\mcite[\( 1.18(\varepsilon) \)]{Ляпин1960Полугруппы} Left and right \term[ru=сокращение, bg=съкращение (\cite[77]{ГеновМиховскиМоллов1991Алгебра}), en=cancellation (\cite[218]{Kleene2002Logic})]{cancellation}:
    \begin{align}
      \qforall \synz (\synz \cdot \synx \syneq \synz \cdot \syny) \synimplies \synx \syneq \syny
      \label{eq:def:binary_operation/cancellative/left}
      \\
      \qforall \synz (\synx \cdot \synz \syneq \syny \cdot \synz) \synimplies \synx \syneq \syny
      \label{eq:def:binary_operation/cancellative/right}
    \end{align}
  \end{thmenum}
\end{definition}

\begin{example}\label{ex:def:binary_operation}
  We list several examples of \hyperref[def:binary_operation]{binary operations} satisfying different properties.

  \begin{thmenum}
    \thmitem{ex:def:binary_operation/algebraic} \hyperref[eq:def:binary_operation/associative]{Associative} \hyperref[def:binary_operation]{binary operations} on a set are abundant and are part of the definition of essential algebraic structures like \hyperref[def:group]{groups}, \hyperref[def:ring]{rings}, \hyperref[def:vector_space]{vector spaces} and \hyperref[def:lattice]{(semi)lattices}.

    These operations are \term{homogeneous} in the sense that their signature only contains a single set, unlike \hyperref[def:group_action]{group actions} and scalar products in \hyperref[def:vector_space]{vector spaces}.

    \thmitem{ex:def:binary_operation/composition} The quintessential example of a non-\hyperref[def:binary_operation/commutative]{commutative} operation is \hyperref[def:set_valued_map/composition]{composition} in any set of functions or, more generally, \hyperref[def:category/composition]{morphism composition} in any \hyperref[def:category]{category}.

    Composition is \hyperref[def:binary_operation/associative]{associative}. \hyperref[def:binary_operation/cancellative]{Cancellation} with respect to composition is discussed in \cref{def:morphism_invertibility} and, for function composition, in \cref{thm:function_invertibility_categorical}.

    \thmitem{ex:def:binary_operation/midpoint} The midpoint operation
    \begin{equation*}
      (x, y) \mapsto \dfrac {x + y} 2
    \end{equation*}
    on \( \BbbR \) is commutative and cancellative, but not associative.

    \thmitem{ex:def:binary_operation/grafting} The \hyperref[def:ordered_tree_grafting_product]{grafting product} \( B_+(T, R) \) can be regarded as a binary operation on \hyperref[def:ordered_tree]{ordered trees}.

    The product of
    \begin{equation*}
      \begin{aligned}
        \includegraphics[page=1]{output/ex__def__binary_operation__grafting}
        \quad
        \includegraphics[page=2]{output/ex__def__binary_operation__grafting}
      \end{aligned}
    \end{equation*}
    is the following tree:
    \begin{equation*}
      \begin{aligned}
        \includegraphics[page=3]{output/ex__def__binary_operation__grafting}
      \end{aligned}
    \end{equation*}

    The product \( B_+ \) clearly depends on the order of the arguments, hence, as a binary operation, it is not commutative.

    It is even simpler to show that \( B_+ \) is not associative --- the following are the trees \( B_+(B_+(\anon, \anon), \anon) \) and \( B_+(\anon, B_+(\anon, \anon)) \), respectively:
    \begin{equation*}
      \begin{aligned}
        \includegraphics[page=4]{output/ex__def__binary_operation__grafting}
        \quad
        \includegraphics[page=5]{output/ex__def__binary_operation__grafting}
      \end{aligned}
    \end{equation*}

    \thmitem{ex:def:binary_operation/collatz}\mcite{MathSE:collatz_as_a_binary_operation} On the \hyperref[def:integer_parity]{odd integers}, consider the operation
    \begin{equation}\label{eq:ex:def:binary_operation/collatz}
      a \ast b \coloneqq \frac {3a + b} {2^{v_2(3a + b)}},
    \end{equation}
    where \( v_2 \) is the \hyperref[def:padic_valuation]{\( 2 \)-adic valuation}, i.e. the largest power of \( 2 \) that divides \( 3a + b \).

    This operation can be used to describe \fullref{hyp:collatz_conjecture}. Indeed, for the \hyperref[def:collatz_map]{Collatz map} \( C \) we have
    \begin{equation*}
      a \ast 1 = \frac {3a + 1} {2^{v_2(3a + 1)}} = C^{v_2(3a + 1) + 1} (a).
    \end{equation*}

    An even integer \( n \) can be reduced to the odd case via division by \( v_2(n) \). Then \fullref{thm:collatz_conjecture_equivalences/recurrence} can be restated as follows:
    \begin{displayquote}
      The integer \( 1 \) is a \hyperref[def:dynamical_system_fixed_point]{fixed point} of the \hyperref[def:recurrence_relation]{recurrence relation} \( X_{N+1} = X_N \ast 1 \) on odd integers. Furthermore, for every starting condition, the recurrence attains this fixed point.
    \end{displayquote}

    As a binary operation, \( a \ast b \) is clearly not commutative because
    \begin{equation*}
      1 \ast 2 = \frac {3 + 2} 1 = 5 \neq 7 = \frac {6 + 1} 1= 2 \ast 1.
    \end{equation*}

    It is also not associative because
    \begin{equation*}
      (1 \ast 2) \ast 3 = 5 \ast 3 = 9 \neq 3 = 1 \ast 9 = 1 \ast (2 \ast 3).
    \end{equation*}
  \end{thmenum}
\end{example}

\begin{concept}\label{con:one_sided_associativity}
  In a \hyperref[def:semigroup]{semigroup}, the products \( ((xy)z) \) and \( (x(yz)) \) are equal due to \hyperref[def:binary_operation/associative]{associativity}. More formally, in the \hyperref[def:fol_free_term_model]{free term model} of the \hyperref[def:semigroup/theory]{first-order theory of semigroups}, the following two syntax trees are congruent:
  \begin{equation*}
    \includegraphics[page=1]{output/con__one_sided_associativity}
    \qquad
    \includegraphics[page=2]{output/con__one_sided_associativity}
  \end{equation*}

  We extensively discuss free term models in \cref{ex:def:fol_free_term_model/free_semigroup}. The key takeaway is that the notation \( xyz \) may refer to either \( ((xy)z) \) or \( (x(yz)) \), but our choice is irrelevant.

  Sometimes we want to use this simpler notation without the operation being associative. Then we have two options:
  \begin{thmenum}
    \thmitem{con:one_sided_associativity/left} We may parse \( xyz \) as \( ((xy)z) \), i.e. as the left syntax tree above. We refer to this convention as \term[en=association to the left (\cite[1]{Hindley1997BasicSTT})]{association to the left}. In this case we leave \( x(yz) \) as-is.

    For example, as discussed in \cref{rem:lambda_term_notation_conventions/left_associative}, it is common to write the \hyperref[def:lambda_term]{\( \muplambda \)-term} \( (MN)K \) as \( MNK \).

    \thmitem{con:one_sided_associativity/right} Alternatively, we may parse \( xyz \) as \( (x(yz)) \), i.e. as the right syntax tree above. We refer to this convention as \term[en=association to the right (\cite[12]{Hindley1997BasicSTT})]{association to the right}. In this case we leave \( (xy)z \) as-is.

    For example, as discussed in \cref{rem:simple_type_parentheses/arrow}, it is common to write the \hyperref[def:arrow_type]{arrow type} \( \tau \synimplies (\sigma \synimplies \rho) \) as \( \tau \synimplies \sigma \synimplies \rho \).
  \end{thmenum}
\end{concept}

\begin{remark}\label{rem:binary_operations_to_infinitary}
  \hyperref[def:binary_operation]{Binary operations} can sometimes be extended to infinite families. To avoid trivialities, we suppose that the operations here are \hyperref[def:binary_operation/associative]{associative} and \hyperref[def:binary_operation/commutative]{commutative} so that the order of the operands does not matter.

  \begin{thmenum}
    \thmitem{rem:binary_operations_to_infinitary/lattice} The operations of a \hyperref[def:lattice]{lattice} arise by specializing \hyperref[def:extremal_points/supremum_and_infimum]{suprema and infima} to binary sets. Conversely, as discussed in \cref{thm:lattice_from_binary_operations}, the binary lattice operations induce a \hyperref[def:partially_ordered_set]{partial order}, which can be used to define arbitrary suprema and infima. If the lattice happens to be \hyperref[def:complete_lattice]{complete}, instead of the binary operations, we can use
    \begin{equation*}
      \bigvee_{k \in \mscrK} x_k = \sup\set{ x_k \given k \in \mscrK }.
    \end{equation*}

    \thmitem{rem:binary_operations_to_infinitary/convergence} Given a sequence \( x_1, x_2, \ldots \) in a \hyperref[def:banach_space]{Banach space}, we can form the corresponding sequence of partial sums \( \sum_{k=1}^n x_k \) and take the limit as \( n \to \infty \).

    This leads to the notion of a \hyperref[def:convergent_series]{series}. Even if the limit does not exist (if \enquote{the series diverges}), the series is still a well-defined syntactic object. This leads us to define formal power series in \cref{def:formal_power_series}.

    \thmitem{rem:binary_operations_to_infinitary/direct_sum} In a \hyperref[def:monoid/commutative]{commutative monoid} \( G \), it is possible to sum infinitely many elements, as long as only finitely many are nonzero.

    Indeed, consider the sum
    \begin{equation*}
      s \coloneqq \sum_{k \in \mscrK} x_k.
    \end{equation*}

    Technically, this must be defined as
    \begin{equation*}
      s \coloneqq \sum_{k \in \mscrK_0} x_k,
    \end{equation*}
    where \( \mscrK_0 \) be the set of indices such that \( x_k \neq 0_M \) when \( k \in \mscrK_0 \).

    It is often more convenient to sum over \( \mscrK \), and we will do is because it is harmless. In fact, this is the reasoning for defining the \hyperref[def:free_commutative_monoid]{free commutative monoid} over \( A \) as the family of finitely-supported functions from \( A \) to \( \BbbN \).
  \end{thmenum}
\end{remark}

\paragraph{Semigroups}

\begin{remark}\label{rem:bourbaki}
  Nicholas Bourbaki is a pseudonym for an evolving group of French mathematicians who have united in the 1930s to produce a systematic exposition of many areas of mathematics. Armand Borel, one of the group's third generation members, writes the following in \cite{Borel1998Bourbaki}:
  \begin{displayquote}
    In the early thirties the situation of mathematics in France at the university and research levels, the only ones of concern here, was highly unsatisfactory. World War I had essentially wiped out one generation. The upcoming young mathematicians had to rely for guidance on the previous one, including the main and illustrious protagonists of the so-called 1900 school, with strong emphasis on analysis. Little information was available about current developments abroad, in particular about the flourishing German school (Göttingen, Hamburg, Berlin), as some young French mathematicians (J. Herbrand, C. Chevalley, A. Weil, J. Leray) were discovering during visits to those centers.

    In 1934 A. Weil and H. Cartan were Ma\^itres de Conf\'erences (the equivalent of assistant professors) at the University of Strasbourg. One main duty was, of course, the teaching of differential and integral calculus. The standard text was the Trait\'e d'Analyse of E. Goursat, which they found wanting in many ways. Cartan was frequently bugging Weil with questions on how to present this material, so that at some point, to get it over with once and for all, Weil suggested they write themselves a new Trait\'e d'Analyse. This suggestion was spread around, and soon a group of about ten mathematicians began to meet regularly to plan this treatise. It was soon decided that the work would be collective, without any acknowledgment of individual contributions. In summer 1935 the pen name Nicolas Bourbaki was chosen.
    \begin{equation*}
      \vdots
    \end{equation*}

    A first question to settle was how to handle references to background material. Most existing books were found unsatisfactory. Even B. v.d. Waerden's \textit{Moderne Algebra}, which had made a deep impression, did not seem well suited to their needs (besides being in German). Moreover, they wanted to adopt a more precise, rigorous style of exposition than had been traditionally used in France, so they decided to start from scratch and, after many discussions, divided this basic material into six \enquote{books}, each consisting possibly of several volumes, namely:
    \begin{enumerate}[label=\Roman*]
      \item Set Theory
      \item Algebra
      \item Topology
      \item Functions of One Real Variable
      \item Topological Vector Spaces
      \item Integration
    \end{enumerate}

    These books were to be linearly ordered: references at a given spot could only be to the previous text in the same book or to an earlier book (in the given ordering). The title \enquote{\'El\'ements de Math\'ematique} was chosen in 1938. It is worth noting that they chose \enquote{Math\'ematique} rather than the much more usual \enquote{Math\'ematiques}. The absence of the \enquote{s} was of course quite intentional, one way for Bourbaki to signal its belief in the unity of mathematics.

    The first volumes to appear were the Fascicle of Results on Set Theory (1939) and then, in the forties, Topology and three volumes of Algebra.
  \end{displayquote}

  The members of the group are recruited among selected guests on their seminars and by tradition retire at 50. A lot of reorganization to the aforementioned planned structure was needed to accommodate both the evolving group and the evolving mathematics. The latest book is \cite{Bourbaki2023ThéoriesSpectrales3à5}, published in 2023. The book \cite{Aluffi2009Algebra}, which attempts to introduce algebra and category theory simultaneously, is a word play on their naming scheme.

  The level of influence of their work varies widely:
  \begin{itemize}
    \item The ubiquitous \hyperref[def:empty_set]{empty set} symbol \( \varnothing \) is attributed to the 1939 edition of \cite[E II.1]{Bourbaki1970ThéorieDesEnsembles} by \cite{Miller2019EarliestUsesOfLogicSymbols}. Another early notation is \( 0 \), used by \incite[12]{Hausdorff1935Mengenlehre} and as late as \bycite[10]{Натансон1974ВещественныйАнализ}.

    \item The ubiquitous \hyperref[def:naive_set_theory]{set membership} symbol \( {\in} \) is also attributed to Bourbaki's 1939 set theory book. Another early notation is \( \varepsilon \), used again by \incite[11]{Hausdorff1935Mengenlehre}, as early as \bycite[X]{Peano1889ArithmeticesPrincipia} and as late as \bycite[1]{Kelley1975GeneralTopology}.

    \item On the other hand, their concept of \hyperref[rem:magma_terminology]{magmas} is not used often, as discussed in \cref{rem:magma_terminology}.

    \item Their convention for capitalizing \hyperref[con:free_construction/indeterminate]{indeterminate} is scarcely used, as discussed in \cref{rem:conventions_for_indeterminates}.
  \end{itemize}

  Borel himself comments on his view of the Bourbaki style of writing before joining the group
  \begin{displayquote}
    I was rather put off by the very dry style, without any concession to the reader, the apparent striving for the utmost generality, the inflexible system of internal references and the total absence of outside ones.
  \end{displayquote}

  \incite[121]{Kuyk1977ComplementarityInMathematics} comments the work of Bourbaki as follows:
  \begin{displayquote}
    A main driving force of this group is the problem of the unity of mathematics. At the beginning of the 20\textsuperscript{th} century, mathematics was a series of disciplines, founded on particular notions, delimited with precision and interconnected in many ways, thereby permitting one area to fructify the other. Continued investigation has shown the existence of a central core, which bridges the many diverse parts. The essence of this unity came to be seen in the progressive systematization of relations existing between the diverse mathematical theories and is today (for the Bourbaki group at least) known under the name of \enquote*{axiomatic method}.

    Now the phrase \enquote*{axiomatic method} is usually an application of formal logic, or deductive reasoning. But this the Bourbaki group considers to be nothing other than an external form which the mathematician gives to his thoughts. It is there for communication and it is useful to have its vocabulary and syntax clearly analysed. But it is still only the surface of mathematics and the least interesting aspect of it at that.

    According to the Bourbaki mathematicians, axiomatics catches what formal logic misses, namely, the profound \textit{intelligibility} of mathematics. Where a superficial observer will see two entirely distinct theories, the mathematical genius may suddenly see a distinct analogy between them in their structure. And it is the task of the axiomatic method to search for such common structures in externally differing theories.
  \end{displayquote}

  Adrian Mathias criticizes Bourbaki's style for diverging from the developments in formal logic, and sticking to obscure formalisms. In \cite[48]{Mathias2013BourbakiCriticism}, he summarizes this as follows:
  \begin{displayquote}
    Hilbert in 1922, in joint work with Bernays, proposed an alternative treatment of predicate logic which, despite its many unsatisfactory aspects, was adopted by Bourbaki for their series of books and by Godement for his treaties on algebra, though leading him to express distrust in logic.

    It is this distrust, intensified to a phobia by the vehemence of Dieudonn\'e's writings, and fostered by, for example, the errors and obscurities of a well-known undergraduate textbook, that has, I suggest, led to the exclusion of logic from the CAPES examination. Centralist rigidity has preserved the underlying confusion and consequently flawed teaching; the recovery will start when mathematicians adopt a post-G\"odelian treatment of logic.
  \end{displayquote}

  The treatment of logic mentioned uses Hilbert's \( \varepsilon \) operator, which we briefly discuss in \cref{con:description_operator/epsilon}.
\end{remark}

\begin{remark}\label{rem:magma_terminology}
  There are two major candidates for a term for a set equipped with a \hyperref[def:binary_operation]{binary operation} but no axioms:
  \begin{itemize}
    \item \enquote{Magma} is used for a set with a binary operation in \hyperref[rem:bourbaki]{Bourbaki}'s \cite[def. A I.1]{Bourbaki1970Algèbre1à3}, and earlier usage can be found the 1954 lectures \cite[18]{Serre1992LieGroups} of Jean-Pierre Serre, one of the group's members. Out of the book we use as references, the only other one using these definitions is \cite[33]{UnivalentFoundationsProgram2013HoTT} (for types).

    When describing the design of the computer algebra system Magma \tcite{in #2, #1}[1]{BosmaCannonMatthews1994MagmaLanguage} explain that the name is based on Bourbaki's definition. Some authors use \enquote{magma} to refer to this system rather than the algebraic structure, for example \incite{Rotman2015AdvancedModernAlgebraPart1}, \incite{ConwaySloane1998SpherePackings} and \incite{CoxLittleOShea2015IdealsVarietiesAlgorithms}.

    \item \enquote{Groupoid} is used for a set with a binary operation in
    \cite[145]{Knauer2019AlgebraicGraphTheory},
    \cite[example 21.3]{Golan1999Semirings} and
    \cite[257]{HillePhillips1996FunctionalAnalysis}. The Russian equivalent \textru{группоид} is used in
    \cite[89]{Мальцев1970АлгебраическиеСистемы},
    \cite[34]{Курош1973ОбщаяАлгебра} and
    \cite[121]{БелоусовТкачёв2004ДискретнаяМатематика}.

    Bourbaki use \enquote{groupo\"ide} in \cite[def. A I.127]{Bourbaki1970Algèbre1à3} for \hyperref[def:binary_operation/cancellative]{cancellative} \hyperref[def:semigroup]{semigroups} with generalized \hyperref[def:monoid_inverse]{inverses}.

    We avoid the term because it coincides with groupoids in category theory, which we define in \cref{def:groupoid}. The term \enquote{groupoid} also refers to the categorical construction in
    \cite[20]{MacLane1998CategoryTheory},
    \cite[def. 1.1.11]{Riehl2016CategoryTheory},
    \cite[12]{Jacobson1989BasicAlgebraII},
    \cite[example 4.6]{Aluffi2009Algebra},
    \cite[145]{Knauer2019AlgebraicGraphTheory} and
    \cite[def. 1.1.29]{Perrone2024StartingCategoryTheory}.

    \item There are also other terms for a set with a binary operation, for example \enquote{\textru{мультипликативное множество}} (\enquote{multiplicative set}), used in \cite[20]{Ляпин1960Полугруппы}.
  \end{itemize}

  We are not interested in the structural theory at such generality, for which reason we will refrain from defining magmas and instead define semigroups in \cref{def:semigroup} as an \hyperref[def:binary_operation/associative]{associative} magma. This latter notion is largely unambiguous --- it is used in
  \cite[144]{MacLane1998CategoryTheory},
  \cite[29]{Jacobson1985BasicAlgebraI},
  \cite[133]{Rotman2015AdvancedModernAlgebraPart1},
  \cite[126]{Aluffi2009Algebra},
  \cite[145]{Knauer2019AlgebraicGraphTheory},
  \cite[def. 8.2.1]{HillePhillips1996FunctionalAnalysis},
  \cite[56]{Savage2008ModelsOfComputation},
  \cite[1]{Golan1999Semirings},
  \cite[2]{Pin1998TropicalSemirings},
  \cite[28]{Ляпин1960Полугруппы},
  \cite[90]{Мальцев1970АлгебраическиеСистемы},
  \cite[34]{Курош1973ОбщаяАлгебра},
  \cite[135]{Кострикин2000АлгебраЧасть1},
  \cite[121]{БелоусовТкачёв2004ДискретнаяМатематика},
  \cite[\S 2.2.1]{Новиков2013ДискретнаяМатематика}.

  Bourbaki define a \enquote{left semi-group} in \cite[exerc. I.2.11]{Bourbaki1998Algebra1to3} as a \hyperref[def:binary_operation/cancellative]{left-cancellative} magma.
\end{remark}

\begin{concept}\label{con:opposite_object}\mimprovised
  We sometimes have a natural notion of \term{opposite object}. We list the relevant for us cases in \cref{rem:con:opposite_object}. In all of them, we have a \hyperref[def:category]{category} \( \cat{C} \) with an \hyperref[def:morphism_invertibility/involution]{involutive} functor \( ({\anon*})^{\oppos} \).
\end{concept}

\begin{remark}\label{rem:con:opposite_object}
  The opposite objects defined in \cref{con:opposite_object}, while nonstandard, simply encompasses the following notions that are more established:
  \begin{itemize}
    \item We have \hyperref[def:semigroup/opposite]{opposite semigroups}, in which the binary operation is reversed, with \hyperref[def:monoid/opposite]{opposite monoids} and \hyperref[def:group/opposite]{opposite groups} as special cases.

    \item We have \hyperref[def:semiring/opposite]{opposite semirings} and \hyperref[def:ring/opposite]{opposite rings}, in which multiplication is reversed.

    \item We have \hyperref[def:preordered_set/opposite]{opposite (pre)ordered sets}, in which the order is reversed, with \hyperref[def:lattice/opposite]{opposite lattices} and \hyperref[def:boolean_algebra/opposite]{opposite Boolean algebras} as special cases.

    These come with important principles --- \fullref{thm:order_duality}, \fullref{thm:lattice_duality} and \fullref{thm:boolean_algebra_duality}.

    \item We have \hyperref[def:directed_multigraph/opposite]{opposite directed multigraphs}, in which the arrows are reversed, with \hyperref[def:directed_graph]{directed graphs} as special cases.

    \item Based on multigraphs, we also have \hyperref[def:opposite_category]{opposite categories} and \hyperref[def:opposite_functor]{opposite functors}.

    Categories also have a duality principle --- \fullref{thm:categorical_principle_of_duality}.
  \end{itemize}
\end{remark}

\begin{definition}\label{def:semigroup}\mcite[28]{Ляпин1960Полугруппы}
  A \term[ru=полугруппа (\cite[28]{Ляпин1960Полугруппы}), en=semigroup (\cite[1]{Golan1999Semirings})]{semigroup} is a nonempty set \( G \) equipped with a \hyperref[def:operation_arity_terminology]{binary} \hyperref[def:operation_on_set]{operation} \( \cdot: G \times G \to G \).

  We often denote this operation by juxtaposition as \( xy \) instead of \( x \cdot y \). Thus written, the operation generalizes \hyperref[def:semiring]{multiplication}, even though we allow the operation to not be commutative in order to also encompass \hyperref[def:set_valued_map/composition]{function composition}. If the operation is meant to generalize addition, we modify the notation and terminology according to \cref{con:additive_semigroup}.

  \begin{thmenum}[series=def:semigroup]
    \thmitem{def:semigroup/theory}\mimprovised We can describe semigroups using \hyperref[def:first_order_logic]{first-order logic}. Over the \hyperref[def:fol_signature]{signature} with one \hyperref[def:fol_signature/notation]{infix} binary \hyperref[def:fol_signature/function]{function symbol}, \( \synmult \), an \hyperref[def:fol_equational_theory]{equational theory} of semigroups is generated by the associativity axiom \eqref{eq:def:binary_operation/associative} for \( \synmult \).

    There is one importance nuance --- we generally allow empty semigroups, however we generally disallow empty \hyperref[def:fol_semantics/model]{empty first-order models}. We discuss this issue in \cref{rem:fol_empty_universe/semantics}.

    \thmitem{def:semigroup/morphism}\mcite[348]{Ляпин1960Полугруппы} Given semigroups \( G \) and \( H \), we say that the function \( \varphi: G \to H \) is a \term[ru=гомоморфизм (полугрупп)]{semigroup homomorphism} if, whenever \( x \) and \( y \) are in \( G \), we have
    \begin{equation}\label{eq:def:semigroup/morphism}
      \varphi(x \cdot_G y) = \varphi(x) \cdot_H \varphi(y).
    \end{equation}

    As long as we allow empty universes for first-order models, these are precisely the \hyperref[def:fol_homomorphism]{first-order homomorphisms}, which makes the concepts from \fullref{sec:first_order_structures} applicable to them.

    \thmitem{def:semigroup/category}\mimprovised For a fixed \hyperref[def:grothendieck_universe]{Grothendieck universe} \( \mscrU \), we can consider the family of all \( \mscrU \)-small semigroups along with their homomorphisms.

    This generalizes the \hyperref[def:fol_theory_model_functor/obj]{category of \( \mscrU \)-small first-order models} by also allowing empty sets.

    \thmitem{def:semigroup/subobject}\mcite[133]{Ляпин1960Полугруппы} A \term[ru=подполугруппа]{sub-semigroup} of \( G \) is a subset \( A \) that is closed under the operation in \( G \), i.e. \( xy \) must belong to \( A \) whenever \( x \) and \( y \) do.

    As long as we allow empty universes for first-order models, \cref{thm:fol_equational_theory_models/substructure} implies that the sub-semigroups of \( G \) are its \hyperref[def:fol_substructure/submodel]{first-order submodels}, and from \cref{thm:fol_categorical_subobjects} that they are its \hyperref[def:categorical_subobject]{categorical subobjects}.

    Furthermore, as a consequence of \cref{thm:fol_equational_theory_models/image}, the \hyperref[def:set_valued_map/image]{image} of a homomorphism \( \varphi: X \to Y \) is a sub-semigroup of \( Y \).

    \thmitem{def:semigroup/opposite}\mcite[def. I.1.2]{Bourbaki1998Algebra1to3} We define the \term{opposite} semigroup \( (G, \cdot)^{\oppos} \) as \( (G, \star) \), where \( {\star} \) reverses the order of multiplicands:
    \begin{equation*}
      x \star y \coloneqq y \cdot x.
    \end{equation*}

    This is an \hyperref[con:opposite_object]{opposite object functor} in categories of semigroups.

    \thmitem{def:semigroup/exponentiation}\mcite[86]{Тыртышников2017ОсновыАлгебры} We define an additional \term[ru=\( n \)-тая степень елемента (группы)]{exponentiation} operation for positive integers \( n \) \hyperref[thm:omega_recursion]{recursively} as
    \begin{equation}\label{eq:def:semigroup/exponentiation}
      x^n \coloneqq \begin{cases}
        x,               &n = 1 \\
        x^{n-1} \cdot x, &n > 1
      \end{cases}
    \end{equation}
  \end{thmenum}
\end{definition}

\begin{definition}\label{def:power_semigroup}\mcite[1]{McCarthyHayes1973PowerSemigroup}
  The \hyperref[def:basic_set_operations/power_set]{power set} \( \pow(G) \) of a \hyperref[def:semigroup]{semigroup} \( (G, \cdot) \) is also a semigroup with the operation
  \begin{equation*}
    A \star B \coloneqq \set{ a \cdot b \given a \in A \T{and} b \in B }.
  \end{equation*}

  We will call this the \term{power semigroup} of \( G \). It comes with the canonical \hyperref[def:fol_embedding]{embedding}
  \begin{equation*}
    \begin{aligned}
      &\iota: G \to \pow(G), \\
      &\iota(g) \coloneqq \set{ g }.
    \end{aligned}
  \end{equation*}
\end{definition}
\begin{comments}
  \item Operations on sets are used for \hyperref[def:topological_vector_space]{topological vector spaces} by \incite[6]{Rudin1991FunctionalAnalysis}, \incite[4]{Rockafellar1997ConvexAnalysis}, \incite[4]{Clarke2013OptimalControl}, \incite[25]{МагарилИльяевТихомиров2002ВыпуклыйАнализ}.

  \item Note that this concept is distinct from \hyperref[def:semiring_ideal/product]{product ideals}, which use the same notation.

  \item Power semigroups have a natural \hyperref[def:ordered_semigroup]{ordering} --- see \cref{ex:def:ordered_semigroup/power}
\end{comments}

\begin{example}\label{ex:def:power_semigroup}
  We list some examples of \hyperref[def:power_semigroup]{power semigroups}:
  \begin{thmenum}
    \thmitem{ex:def:power_semigroup/cancellative} The cancellative property does not extend to the power semigroup.

    Consider the additive group \hyperref[def:group_of_integers_modulo]{\( \BbbZ_2 \)}, which is is cancellative as a consequence of \cref{thm:def:group/cancellative}. Define the sets \( A \coloneqq \set{ 0, 1 } \) and \( B \coloneqq \set{ 0 } \). Then
    \begin{equation*}
      A + A = A = A + B,
    \end{equation*}
    however we cannot cancel \( A \) from the left because \( A \neq B \).
  \end{thmenum}
\end{example}

\paragraph{Semigroup exponentiation}

\begin{definition}\label{def:semigroup_power}\mimprovised
  In the context of \hyperref[def:semigroup/exponentiation]{semigroup exponentiation}, if \( y = x^n \), we say that \( y \) is the \term[ru=\( n \)-тая степень (\cite[20]{Ляпин1960Полугруппы})]{\( n \)-th power} of \( x \) if \( x^n = y \) and that \( x \) is an \term{\( n \)-th root} of \( y \). If \( y \) has an \( n \)-th root, we call it a \term{perfect \( n \)-th power}.

  We will use the following terminology based on \hyperref[def:hypercubic_number]{hypercubic numbers}:
  \begin{center}
    \begin{tabular}{lll}
      \toprule
      \( n \) & \( x \)            & \( y \)       \\
      \midrule
      \( 2 \) & \term{square root} & \term{square} \\
      \( 3 \) & \term{cubic root}  & \term{cube}   \\
      \bottomrule
    \end{tabular}
  \end{center}
\end{definition}
\begin{comments}
  \item While \( y \) is the unique \( n \)-th power of \( x \), \( x \) is only one of many possible roots of \( y \).

  \item The term \enquote{root} is discussed in \cref{rem:root_terminology}, along with how it relates to \hyperref[def:polynomial_root]{polynomial roots}, \hyperref[def:polynomial_equation/solution]{solutions} to \hyperref[def:polynomial_equation]{algebraic equations} and \hyperref[def:zero_of_function]{zeros of functions}.

  \item In the case of positive \hyperref[def:real_numbers]{real numbers}, we have a canonical choice of \( n \)-th roots given by \cref{def:principal_nonnegative_nth_root}, for which we use the notation \( \sqrt[n]{ x } \). We also have such a canonical choice for square roots of negative real numbers --- see \cref{def:principal_real_square_root}.
\end{comments}

\begin{remark}\label{rem:square_and_quadratic}
  The \hyperref[def:semigroup_power]{square root} is a solution to a \hyperref[tab:def:polynomial_degree]{quadratic} \hyperref[def:polynomial_equation]{equation}, although a \hyperref[def:semigroup_power]{cubic root} is a solution to a \hyperref[tab:def:polynomial_degree]{cubic} equation.

  According to \incite[8]{ConwayGuy1998BookOfNumbers}, both \enquote{square} and \enquote{quadratic} derive from the Latin \enquote{\foreignlanguage{latin}{ex quadrem}}. In Russian, \enquote{square root} and \enquote{quadratic equation} both use the same adjective, \enquote{\textru{квадратный}} (\enquote{quadratic}), however the established English terms are for some reason different.
\end{remark}

\begin{proposition}\label{thm:semigroup_exponentiation_properties}
  \hyperref[def:semigroup/exponentiation]{Exponentiation} in a fixed \hyperref[def:semigroup]{semigroup} \( G \) has the following basic properties:

  \begin{thmenum}
    \thmitem{thm:semigroup_exponentiation_properties/commutativity} We have the following \hyperref[def:binary_operation/commutative]{commutativity}-like property: for any \( x \) in \( G \) and any positive integer \( n \), we have
    \begin{equation}\label{eq:thm:semigroup_exponentiation_properties/commutativity}
      x^n = x x^{n-1} = x^{n-1} x.
    \end{equation}

    \thmitem{thm:semigroup_exponentiation_properties/multiplication} If \( G \) is \hyperref[def:semiring/commutative]{commutative}, then, for any two members \( x \) and \( y \) of \( G \) and any positive integer \( n \),
    \begin{equation}\label{eq:thm:semigroup_exponentiation_properties/multiplication}
      (xy)^n = x^n \cdot y^n
    \end{equation}

    \thmitem{thm:semigroup_exponentiation_properties/distributive} Exponentiation distributes over multiplication: for any member \( x \) of \( G \) and any two positive integers \( n \) and \( m \),
    \begin{equation}\label{eq:thm:semigroup_exponentiation_properties/distributive}
      x^{n + m} = x^n x^m.
    \end{equation}

    \thmitem{thm:semigroup_exponentiation_properties/repeated} For any member \( x \) of \( M \) and any two positive integers \( n \) and \( m \),
    \begin{equation}\label{eq:thm:semigroup_exponentiation_properties/repeated}
      (x^n)^m = x^{nm}.
    \end{equation}
  \end{thmenum}
\end{proposition}
\begin{proof}
  \SubProofOf{thm:semigroup_exponentiation_properties/commutativity} We use induction on \( n \). The cases \( n = 1 \) and \( n = 2 \) are obvious. For \( n > 2 \), we have
  \begin{equation*}
    x^n
    \reloset {\eqref{eq:def:semigroup/exponentiation}} =
    x x^{n-1}
    \reloset {\T{ind.}} =
    x x^{n-2} x
    \reloset {\eqref{eq:def:semigroup/exponentiation}} =
    x^{n-1} x.
  \end{equation*}

  \SubProofOf{thm:semigroup_exponentiation_properties/multiplication} We use induction on \( n \). The case \( n = 1 \) is obvious. For \( n > 1 \), we have
  \begin{equation*}
    (xy)^n = x^n \cdot y^n
    \reloset {\eqref{eq:def:semigroup/exponentiation}} =
    (xy) \cdot (xy)^{n-1}
    \reloset {\T{ind.}} =
    (xy) \cdot x^{n-1} \cdot y^{n-1}
    \reloset {\eqref{eq:def:binary_operation/commutative}} =
    x \cdot x^{n-1} \cdot y \cdot y^{n-1}
    =
    x^n \cdot y^n.
  \end{equation*}

  \SubProofOf{thm:semigroup_exponentiation_properties/distributive} We use induction on \( n \). The case \( n = 1 \) follows directly from \eqref{eq:def:semigroup/exponentiation}. The case \( n > 1 \) follows from
  \begin{equation*}
    x^{n + m}
    \reloset {\eqref{eq:def:semigroup/exponentiation}} =
    x x^{n + (m - 1)}
    \reloset {\T{ind.}} =
    x x^{n-1} x^m
    \reloset {\eqref{eq:def:semigroup/exponentiation}} =
    x^n x^m.
  \end{equation*}

  \SubProofOf{thm:semigroup_exponentiation_properties/repeated} We use induction on \( n \). The case \( n = 1 \) is obvious and the rest follows from
  \begin{equation*}
    (x^n)^m
    \reloset {\eqref{eq:def:semigroup/exponentiation}} =
    x^n (x^n)^{m-1}
    \reloset {\T{ind.}} =
    x^n x^{n (m - 1)}
    \reloset {\eqref{eq:thm:semigroup_exponentiation_properties/multiplication}} =
    =
    x^{nm}.
  \end{equation*}
\end{proof}

\paragraph{Ordered semirings}

\begin{definition}\label{def:ordered_semigroup}\mcite[535]{Ляпин1960Полугруппы}
  A (partially) \term[ru=частично упорядоченная полугруппа (\cite[535]{Ляпин1960Полугруппы})]{ordered semigroup} is a \hyperref[def:semigroup]{semigroup} \( M \) equipped with a \hyperref[def:partially_ordered_set]{partial order} \( \leq \) such that \( x \leq y \) implies \( xz \leq yz \) and \( zx \leq zy \) for every \( z \) in \( M \).
\end{definition}
\begin{comments}
  \item The category of small ordered semigroups is \hyperref[def:concrete_category]{concrete} over both the \hyperref[def:semigroup/category]{category of semigroups} and the \hyperref[def:partially_ordered_set]{category of partially ordered sets}.
\end{comments}

\begin{proposition}\label{thm:def:ordered_semigroup}
  \hyperref[def:ordered_semigroup]{Ordered semigroups} have the following basic properties:
  \begin{thmenum}
    \thmitem{thm:def:ordered_semigroup/cancellative} If the semigroup is \hyperref[def:binary_operation/cancellative]{cancellative}, then \( x < y \) implies \( xz < yz \) and \( zx < zy \).
  \end{thmenum}
\end{proposition}
\begin{comments}
  \item \hyperref[def:group]{Groups} behave better with respect to ordering --- see \cref{thm:ordered_group}.
\end{comments}
\begin{proof}
  \SubProofOf{thm:def:ordered_semigroup/cancellative} Trivial.
\end{proof}

\begin{example}\label{ex:def:ordered_semigroup}
  We list several examples of \hyperref[def:ordered_semigroup]{ordered semigroups}:

  \begin{thmenum}
    \thmitem{ex:def:ordered_semigroup/natural_numbers} The \hyperref[def:natural_numbers]{natural numbers} with addition form an ordered semigroup as a consequence of \cref{thm:natural_numbers_are_well_ordered}; and so do \( \BbbZ \), \( \BbbQ \), \( \BbbR \) and \( \BbbC \).

    \thmitem{ex:def:ordered_semigroup/ordinal} More generally, every \hyperref[def:successor_and_limit_ordinal]{limit ordinal} as the set of all smaller ordinals is an ordered semigroup under \hyperref[def:ordinal_arithmetic/addition]{ordinal addition}.

    Unlike addition in natural numbers, however, ordinal addition is not commutative, as shown in \cref{ex:ordinal_addition}.

    \thmitem{ex:def:ordered_semigroup/semilattice} Every \hyperref[def:semilattice]{join-semilattice} \( (X, \vee) \) is an ordered semigroup with the lattice order. Indeed, if \( x \leq y \), then
    \begin{itemize}
      \item If \( z \leq x \), then
      \begin{equation*}
        x = x \vee z \leq y \vee z = y.
      \end{equation*}

      \item If \( x \leq z \leq y \), then
      \begin{equation*}
        z = x \vee z \leq y \vee z = y.
      \end{equation*}

      \item If \( z \geq y \), then
      \begin{equation*}
        z = x \vee z \leq y \vee z = z.
      \end{equation*}
    \end{itemize}

    Every \hyperref[def:semilattice]{meet-semilattice} is similarly an ordered semigroup.

    \thmitem{ex:def:ordered_semigroup/power} Any \hyperref[def:power_semigroup]{power semigroup} is a Boolean algebra as discussed in \cref{thm:boolean_algebra_of_subsets}. As shown in \cref{ex:def:ordered_semigroup/semilattice}, this makes it an ordered semigroup.
  \end{thmenum}
\end{example}

\paragraph{Monoids}

\begin{definition}\label{def:monoid}\mcite[def. 1.1]{Jacobson1985BasicAlgebraI}
  A \term[ru=моноид (\cite[94]{Мальцев1970АлгебраическиеСистемы})]{monoid} is a \hyperref[def:semigroup]{semigroup} with a distinguished element \( e \) such that \( xe = x = ex \) for every \( x \). Such an element is obviously unique, and we call it the \term[ru=нейтральный елемент (\cite[115]{БелоусовТкачёв2004ДискретнаяМатематика}), en=neutral element (\cite[165]{Knauer2019AlgebraicGraphTheory})]{neutral element} of the monoid. This makes monoids \hyperref[def:pointed_set]{pointed sets}.

  \begin{thmenum}
    \thmitem{def:monoid/theory}\mimprovised For monoids, we extend the \hyperref[def:semigroup/theory]{first-order signature for semigroups} with a nullary function symbol, \( \syn1 \), based on the \hyperref[def:pointed_set]{signature for pointed sets}.

    To obtain an \hyperref[def:fol_equational_theory]{equational theory} of monoids over this signature, we extend the theory of semigroups with the axiom
    \begin{align}
      \qforall \synx (\synx \synmult \syn1 \syneq \synx), \label{eq:def:monoid/theory/neutral/left} \\
      \qforall \synx (\syn1 \synmult \synx \syneq \synx). \label{eq:def:monoid/theory/neutral/right}
    \end{align}

    \thmitem{def:monoid/morphism}\mcite[def. 1.6]{Jacobson1985BasicAlgebraI} Given monoids \( M \) and \( N \), we say that the function \( \varphi: M \to N \) is a \term[en=homomorphism (of monoids) (\cite[def. 1.6]{Jacobson1985BasicAlgebraI})]{monoid homomorphism} if it is both a \hyperref[def:semigroup/morphism]{semigroup homomorphism} and a \hyperref[def:pointed_set/morphism]{pointed set homomorphism}, i.e. it respects both the monoid operation and the neutral element.

    These are precisely the \hyperref[def:fol_homomorphism]{first-order homomorphisms}, which makes the concepts from \fullref{sec:first_order_structures} applicable to them.

    \thmitem{def:monoid/category}\mcite[12]{MacLane1998CategoryTheory} For a fixed \hyperref[def:grothendieck_universe]{Grothendieck universe} \( \mscrU \), the first-order theory of monoids naturally gives rise to a \hyperref[def:fol_theory_model_functor/obj]{category of \( \mscrU \)-small models}.

    We denote this category by \( \ucat{Mon} \).

    We list several properties:
    \begin{itemize}
      \item \( \ucat{Mon} \) is \hyperref[def:concrete_category]{concrete} over both the \hyperref[def:pointed_set/category]{category of \( \mscrU \)-small pointed sets} and the \hyperref[def:pointed_set/category]{category of \( \mscrU \)-small semigroups}.

      \item As an equational theory, the theory of monoids has \hyperref[def:fol_free_term_model]{free term models}, and \fullref{thm:fol_free_term_model_universal_property} provides a \hyperref[def:category_adjunction]{left adjoint} to the \hyperref[def:fol_theory_forgetful_functor/set]{forgetful functor} \( U: \ucat{Mon} \to \ucat{Set} \).

      A more convenient construction is the \hyperref[def:free_monoid]{free monoid}, for which we have \fullref{thm:free_monoid_universal_property}.

      The two constructions are necessarily isomorphic because of \cref{thm:functor_adjoint_uniqueness}.

      \item Because of this left adjoint, \cref{thm:concrete_category_function_invertibility/injective} implies that a monoid homomorphism is a \hyperref[def:morphism_invertibility/left_cancellative]{monomorphism} if and only if it is \hyperref[def:function_invertibility/injective]{injective}.

      \item For epimorphisms a similar statement does not hold, unfortunately. The embedding \( \iota: \BbbN \to \BbbZ \) is an epimorphism due to \fullref{thm:grothendieck_monoid_completion_universal_property}, but it is clearly not surjective.
    \end{itemize}

    \thmitem{def:monoid/subobject}\mcite[29]{Jacobson1985BasicAlgebraI} A \term[ru=подмоноид (\cite[148]{БелоусовТкачёв2004ДискретнаяМатематика})]{submonoid} of \( M \) is a \hyperref[def:semigroup/subobject]{sub-semigroup} that contains the neutral element \( e \).

    It follows from \cref{thm:fol_equational_theory_models/substructure} that the submonoids of \( M \) are precisely its \hyperref[def:fol_substructure/submodel]{first-order submodels}, and from \cref{thm:fol_categorical_subobjects} that they are its \hyperref[def:categorical_subobject]{categorical subobjects}.

    Furthermore, as a consequence of \cref{thm:fol_equational_theory_models/image}, the image of a monoid homomorphism is a submonoid of its codomain.

    \thmitem{def:monoid/trivial}\mcite[exerc. II.1.3]{Aluffi2009Algebra} In accordance with \cref{def:trivial_object}, we will call a one-element monoid \term{trivial} because, as shown in \cref{thm:equational_theory_zero_morphisms/trivial}, it is a zero object in \( \ucat{Mon} \).

    \thmitem{def:monoid/opposite}\mimprovised The \hyperref[def:semigroup/opposite]{opposite semigroup} \( M^{\oppos} \) of a monoid \( M \) is a monoid with the same neutral element.

    For this reason, we will call it the \term{opposite monoid}.

    \thmitem{def:monoid/exponentiation}\mcite[86]{Тыртышников2017ОсновыАлгебры} We extend \hyperref[def:semigroup/exponentiation]{semigroup exponentiation} to all nonnegative integers by additionally defining
    \begin{equation*}
      x^0 \coloneqq e.
    \end{equation*}

    \thmitem{def:monoid/commutative}\mcite[41]{Jacobson1985BasicAlgebraI} If the monoid operation is \hyperref[def:binary_operation/commutative]{commutative}, we say that the monoid itself is \term{commutative}.
  \end{thmenum}
\end{definition}

\begin{example}\label{ex:def:monoid}
  We list examples of \hyperref[def:monoid]{monoids}.

  \begin{thmenum}
    \thmitem{ex:def:monoid/natural_numbers} The \hyperref[def:natural_numbers]{nonnegative natural numbers} with addition form a quintessential example of a monoid. We prove in \cref{thm:natural_number_addition_properties} that \( (\BbbN, +) \) is indeed a monoid.

    \thmitem{ex:def:monoid/weak_limit_cardinal} More generally, for every \hyperref[def:successor_and_limit_cardinal/weak_limit]{weak limit cardinal}, the set of smaller cardinals is a commutative monoid under \hyperref[def:cardinal_arithmetic/addition]{cardinal addition} as a consequence of \cref{thm:cardinal_addition_algebraic_properties}.

    \thmitem{ex:def:monoid/limit_ordinal} For a \hyperref[def:successor_and_limit_ordinal]{limit ordinal}, however, the set of smaller ordinals is \hi{non-commutative} under \hyperref[def:ordinal_arithmetic/addition]{ordinal addition} as a consequence of \cref{thm:ordinal_addition_algebraic_properties}.

    \thmitem{ex:def:monoid/kleene_star} Another important example of a monoid is the \hyperref[def:formal_language/kleene_star]{Kleene star} \( \mscrA \) over some \hyperref[def:formal_language]{alphabet} \( \mscrA \).

    The importance for monoid theory comes from the free monoid universal property described in \fullref{thm:free_monoid_universal_property}.

    \thmitem{ex:def:monoid/semilattice} If a \hyperref[def:semilattice]{join-semilattice} is bounded from below, it is a monoid. A similar statement holds for \hyperref[def:lattice/meet]{meet-semilattices} that are bounded from above.

    \thmitem{ex:def:monoid/power} The \hyperref[def:power_semigroup]{power semigroup} \( \pow(M) \) of a monoid \( M \) is also a monoid --- \( \iota(e_M) = \set{ e_M } \) is a neutral element.
  \end{thmenum}
\end{example}

\begin{example}\label{ex:monoid_cancellation_not_preserved_by_homomorphism}\mcite{MathSE:semigroup_cancellation_not_preserved}
  \hyperref[def:monoid/morphism]{Monoid homomorphisms} may not preserve the \hyperref[def:binary_operation/cancellative]{cancellation property}. For example, the \hyperref[def:natural_numbers]{natural numbers} \( \BbbN \) are a cancellative monoid under addition, as shown in \cref{thm:natural_number_addition_properties}, but the homomorphism
  \begin{equation*}
    \begin{aligned}
      &h: (\BbbN, +) \to (\BbbF_2, \max) \\
      &h(n) \coloneqq \begin{cases}
        0, &n = 0 \\
        1, &n > 0
      \end{cases}
    \end{aligned}
  \end{equation*}
  does not preserve the cancellation property.

  Indeed, \( \max\set{ 0, 1 } = \max\set{ 1, 1 } \), but \( 0 \neq 1 \).
\end{example}

\begin{definition}\label{def:monoid_idempotent}\mcite[1]{Golan1999Semirings}
  We say that an element \( x \) of a \hyperref[def:monoid]{monoid} (or, more generally, a \hyperref[def:semigroup]{semigroup}) is \term[ru=идемпотент (\cite[72]{Ляпин1960Полугруппы}), en=idempotent]{idempotent} if \( x^2 = x \). Since the neutral element is always idempotent, we call idempotent elements distinct from it \term[en=nontrivial (idempotent) (\cite[287]{Rudin1991FunctionalAnalysis})]{nontrivial}.
\end{definition}
\begin{comments}
  \item The monoid operation is itself idempotent in the sense of \cref{def:binary_operation/idempotent} if and only if every element is idempotent.
\end{comments}
